\documentclass[11pt]{article}
\usepackage[margin=1in]{geometry}
\usepackage{amsmath,amssymb,amsthm,enumitem,hyperref}
\hypersetup{colorlinks=true,linkcolor=blue,urlcolor=blue,citecolor=blue}
\newtheorem{theorem}{Theorem}
\newcommand{\HH}{\mathcal H}
\title{Infinitely many failures of binary Hankel monotonicity\\\large MathDB 375936}
\author{Principia Math In Progress}
\date{September 18, 2026}
\begin{document}
\maketitle
\begin{abstract}
Both monotonicity assertions in Sobolewski--Ulas's Conjecture 4.7
are false. The adjacent values at indices $25,26$ contradict the
increasing assertion, and those at $39,40$ contradict the decreasing
assertion. We use the source's recurrence to propagate these
violations to infinitely many intervals, with explicit signed formulas.
The small determinant values already occur in OEIS A056886.
\end{abstract}

\section*{Task 1. Full disproof}
Let $s(a)$ count the ones in the binary expansion of $a$, and define
\[
 H(n)=\det[s(i+j)]_{0\le i,j<n},\qquad
 \nu_k=\left\lceil\frac{2^{k+2}}3\right\rceil.
\]
Thus $H(n)=\HH(n,1)$ in \cite{SU}. Conjecture 4.7 asserts that
$|H(n)|$ is strictly increasing on $[\nu_k,2^{k+1}+1]$ and
strictly decreasing on $[2^{k+1}+1,\nu_{k+1}]$, for every $k\ge0$.
Already
\[
 H(25)=-738,\quad H(26)=351,\quad
 H(39)=-432,\quad H(40)=-927                              \tag{1}
\]
give the required counterexamples. Indeed, $25<26$ lie inside
$[\nu_4,33]=[22,33]$, but $738>351$. Likewise $39<40$ lie inside
$[33,\nu_5]=[33,43]$, but $432<927$.
The values (1) are independently certified below and are also
tabulated in \cite{OEIS}; we do not claim their numerical discovery.

\begin{theorem}
For every $k\ge4$, set
\[
 a_k=\nu_k+3(-1)^k,\qquad b_k=\nu_k+4(-1)^k.
\]
Then
\[
 \begin{aligned}
 H(a_k)&=(-1)^{\binom{k}{2}+1}(189k-18),\\
 H(b_k)&=(-1)^{\binom{k+1}{2}}(81k+27).
 \end{aligned}                                               \tag{2}
\]
Every even $k\ge4$ gives a violation in the increasing interval
beginning at $\nu_k$. Every odd $k\ge5$ gives a violation in the
decreasing interval ending at $\nu_k$.
\end{theorem}

\subsection*{Exact seed certificates}
For $2^q<n\le2^{q+1}$, let $h(n,x)$ be the weighted digit-sum
Hankel determinant with weights $1$ at positions $0,\ldots,q$ and
$x$ at position $q+1$. Let $g(n,x)$ be the determinant whose first
$n-1$ columns contain the first differences of that digit-sum
sequence and whose last column is $1$. Thus $H(n)=h(n,1)$.
Put $G(n)=g(n,1)$.
The established source recurrence \cite[Proposition 4.3]{SU}, at $t=1$,
reads
\[
 \begin{pmatrix}h(n,x)\\g(n,x)\end{pmatrix}
 =\begin{cases}
 (-1)^n(x-2)^{2n-2^{q+1}-1}
 \begin{pmatrix}-1&\dfrac{(-1)^n}{x-2}\\0&1\end{pmatrix}
 \begin{pmatrix}h(2^{q+1}-n+1,1)\\g(2^{q+1}-n+1,1)\end{pmatrix},
 &n\le3\cdot2^{q-1},\\[6pt]
 2^{2n-3\cdot2^q-1}(x-2)^{2^q}
 \begin{pmatrix}1&\dfrac{(-1)^n}{x-2}\\0&1\end{pmatrix}
 \begin{pmatrix}h(n-2^q,x/2)\\g(n-2^q,x/2)\end{pmatrix},
 &n>3\cdot2^{q-1}.
 \end{cases}                                                  \tag{3}
\]
Here $n\ge5$. Direct determinants of order at most $4$ give
\[
 \begin{aligned}
 h(3,x)&=3-x,&g(3,x)&=2-x,\\
 h(4,x)&=(x-2)(3x-8),&g(4,x)&=(x-2)^2(x-4).
 \end{aligned}
\]
These bases and (3) give the following exact table:
\[
 \begin{array}{c|c|r|r}
 n&x&h(n,x)&g(n,x)\\\hline
 3&1/2&5/2&3/2\\
 4&1/2&39/4&-63/8\\
 7&1&8&3\\
 8&1&141&-63\\
 9&1/2&-549/2&-189/2\\
 10&1/2&135/4&-81/8\\
 25&1&-738&-189\\
 26&1&351&-81
 \end{array}                                                  \tag{4}
\]
Each row after the first two uses only earlier rows. This proves
the first two values in (1) by exact rational arithmetic.

\subsection*{Propagation to infinitely many violations}
Specializing the first branch of (3) to $x=1$ gives, for
$2^q<N\le3\cdot2^{q-1}$ and $n=2^{q+1}-N+1$,
\[
 H(N)=(-1)^NH(n)+G(n),\qquad
 G(N)=(-1)^{N+1}G(n).                                        \tag{5}
\]
Suppose $H(n),G(n)\ne0$ and
\[
 \operatorname{sgn}G(n)=(-1)^{n+1}\operatorname{sgn}H(n).        \tag{6}
\]
Since $N+n$ is odd, the two terms for $H(N)$ in (5) have the
same sign. Therefore
\[
 |H(N)|=|H(n)|+|G(n)|,\qquad |G(N)|=|G(n)|,                   \tag{7}
\]
and (6) holds at $N$ as well. Condition (6) holds for both seeds
$n=25,26$ in (4).

The source's elementary identity
$\nu_{k+1}=2^{k+2}-\nu_k+1$ gives
\[
 a_{k+1}=2^{k+2}-a_k+1,\qquad
 b_{k+1}=2^{k+2}-b_k+1.                                     \tag{8}
\]
Both new indices lie in the first-branch range for $q=k+1$.
For completeness, with $Q=2^{k+1}\ge32$, each new index lies
between $4Q/3-4$ and $4Q/3+5$, hence strictly between $Q$ and
$3Q/2$. Thus (5)--(7) apply at every step.
Starting from (4), they prove
\[
 |H(a_k)|=738+189(k-4)=189k-18,\qquad
 |H(b_k)|=351+81(k-4)=81k+27.                                \tag{9}
\]
The parity relation $\nu_k\equiv k\pmod2$, obtained from (8),
and the signs in (5) give the signs displayed in (2).

The difference of the magnitudes in (9) is $108k-45>0$.
If $k\ge4$ is even, $a_k<b_k$ lie in
$[\nu_k,2^{k+1}+1]$, contradicting increase. If $k\ge5$ is
odd, $b_k<a_k$ lie in $[2^k+1,\nu_k]$, contradicting decrease.
These interval inclusions follow directly from the formula for
$\nu_k$; the distances $3,4$ fit for the stated bounds on $k$.
At $k=5$, (2) also gives the last two values in (1).
This disproves both universal monotonicity assertions.

The separate record-maximum statement mentioned after the conjecture
is weaker and is not decided by these counterexamples.

\section*{Verification and attribution}
The verifier independently computes all four determinants in (1)
from their defining integer matrices using two exact algorithms.
It also checks the auxiliary seed determinants, the rational table,
and the signed formulas and interval inclusions at larger indices.
The proof of the infinite family is (5)--(9), not finite testing.
The source recurrence belongs to Sobolewski and Ulas, and the small
values in (1) are prior OEIS data. The contribution here is their
explicit application against Conjecture 4.7 and the infinite families.
Review is by the authoring agent only.

\section*{Tasks 2--5. Follow-up directions}
\begin{enumerate}[start=2,leftmargin=*]
\item Decide the weaker record-maximum assertion
$|H(2^k+1)|=\max_{n\le2^k+1}|H(n)|$.
\item Classify all adjacent increases and decreases in $|H(n)|$
using the binary expansion of $n$ and the source recurrence.
\item For every fixed offset $d$, determine eventual signed formulas
for $H(\nu_k+(-1)^kd)$ and the corresponding auxiliary determinants.
\item Quantify the number and density of monotonicity violations in
each interval originally proposed in Conjecture 4.7.
\end{enumerate}

\begin{thebibliography}{9}
\raggedright
\bibitem{SU} B. Sobolewski and M. Ulas,
\emph{Hankel determinants of weighted binary sums of digits},
\href{https://arxiv.org/abs/2607.09376v1}{arXiv:2607.09376v1}, July 10, 2026.
Proposition 4.3 and Conjecture 4.7.
\bibitem{OEIS} OEIS Foundation, \href{https://oeis.org/A056886}{A056886}.
Entry by J. Shallit; expanded table by A. P. Heinz.
\bibitem{DB} MathDB, \href{https://mathdb.com/p/375936}{Problem 375936},
accessed September 18, 2026.
\end{thebibliography}
\end{document}
